\chapter{Conclusion and Future Work}\label{ch:concl}

\section{Strengths and Limitations}

\textbf{Strengths.} The system reaches competitive accuracy with modest
compute. Subword alignment handles agglutination without a morphological
analyser. The CRF produces structurally valid spans. The evaluation uses
strict span-level metrics, and the code is reproducible and deployable.

\textbf{Limitations.}
\begin{itemize}
  \item Only 50k of the 716k training sentences were used, from a single
  run with no hyper-parameter search.
  \item The softmax and CRF models differ in both decoder and loss
  weighting, so the two effects are not separated.
  \item \texttt{pytorch-crf} assumes a contiguous mask. Our training mask
  skips continuation subwords while decoding covers all subwords, so training
  and inference see slightly different paths. This explains the 20 invalid
  transitions that remain.
  \item The training labels are projected from English and are noisy,
  especially for ORG boundaries.
  \item DATE/CARDINAL rules are not evaluated, and the displayed confidence
  scores are not calibrated.
\end{itemize}

\section{Future Work}

\begin{itemize}
  \item Apply the CRF over word-level (first-subword) emissions with hard
  IOB2 constraints, which guarantees valid output.
  \item Run ablations (CRF with weights, softmax without weights) over
  several seeds.
  \item Benchmark mBERT, MuRIL and IndicBERTv2 under identical settings, and
  train on the full corpus.
  \item Improve ORG recognition with gazetteers or span-based models.
  \item Deploy on CPU with ONNX and INT8 quantisation, and add document
  highlighting and entity linking to Malayalam Wikipedia.
\end{itemize}

\section{Conclusion}

We fine-tuned the multilingual IndicBERTv2 encoder for Malayalam NER on
Naamapadam, using first-subword label alignment and a CRF decoder. The model
reaches a strict entity-level F1 of \textbf{84.25\%} on the human-annotated
test set. This is 12 points above a softmax baseline and higher than the
Malayalam score reported in the base paper. The main remaining challenge is
multi-word organization names. Packaged as a web app and a REST API, the
system shows that multilingual Transformers can bring practical information
extraction to Malayalam with modest resources.
